%------------------------------------------------------------------------------%
\begin{question}
  Ortonormalize a base
  \[
    B=
    \{
    (1,1,1),
    (1,0,1),
    (0,1,1)
    \}
  \]
  de $\R^3$ pelo processo de Gram--Schmidt
\end{question}

%------------------------------------------------------------------------------%
\begin{resolution}{Solução}
  \[
    \vec{u}_1
    \pause
    = (1, 1, 1)
  \]

  \pause
  \[
    \norm{\vec{u}_1}
    \pause
    = \sqrt{3}
  \]

  \pause
  \[
    \vec{e}_1
    \pause
    = \frac{1}{\sqrt{3}}(1,1,1)
  \]
\end{resolution}

%------------------------------------------------------------------------------%
\begin{resolution}{Solução}
\begin{columns}[t]
\column{0.4\textwidth}
  \begin{align*}
    \onslide<+->{\vec{u}_2}
    \onslide<+->{& = (1,0,1) - \frac{(1,0,1)\cdot(1,1,1)}{3}\, (1,1,1) \\}
    \onslide<+->{& = (1,0,1) - \frac{2}{3}\, (1,1,1) \\}
    \onslide<+->{& = \left(
        \frac{1}{3},
        -\frac{2}{3},
        \frac{1}{3}
      \right) \\}
    \onslide<+->{& = \frac{1}{3}(1,-2,1)}
  \end{align*}
\column{0.4\textwidth}
  \begin{align*}
    \onslide<+->{\norm{\vec{u}_2}}
    \onslide<+->{& = \frac{\sqrt{6}}{3}}
  \end{align*}
  \begin{align*}
    \onslide<+->{\vec{e}_2}
    \onslide<+->{& = \frac{\vec{u}_2}{\norm{\vec{u}_2}}\\}
    \onslide<+->{& = \frac{1}{\sqrt{6}}(1,-2,1)}
  \end{align*}
\end{columns}
\end{resolution}

%------------------------------------------------------------------------------%
\begin{resolution}{Solução}
\begin{align*}
  \onslide<+->{\vec{u}_3}
  \onslide<+->{& = (0,1,1)
    - \frac{2}{3}(1,1,1)
    -\frac{ (0,1,1)\cdot(1,-2,1) }{6} (1,-2,1)}
  \\
  \onslide<+->{& = (0,1,1)
    - \frac{2}{3}(1,1,1)
    + \frac{1}{6}(1,-2,1)}
  \\
  \onslide<+->{& = \left(
        -\frac{1}{2},
        0,
        \frac{1}{2}
      \right)}
  \\
  \onslide<+->{& = \frac{1}{2}(-1, 0, 1)}
\end{align*}
\end{resolution}

%------------------------------------------------------------------------------%
\begin{resolution}{Solução}
  \begin{align*}
    \onslide<+->{\norm{\vec{u}_3}}
    \onslide<+->{& = \frac{\sqrt{2}}{2}}
  \end{align*}
  \begin{align*}
    \onslide<+->{\vec{e}_3}
    \onslide<+->{& = \frac{\vec{u}_3}{\norm{\vec{u}_3}} \\}
    \onslide<+->{& = \frac{1}{\sqrt{2}}(-1, 0, 1)}
  \end{align*}

  \onslide<+->{\[
    \mathcal{B}
    = \left\{
      \frac{1}{\sqrt{3}}(1,1,1), \;
      \frac{1}{\sqrt{6}}(1,-2,1), \;
      \frac{1}{\sqrt{2}}(-1,0,1)
    \right\}
  \]}
\end{resolution}

%%------------------------------------------------------------------------------%
%\begin{resolution}{Verificação}
%\begin{align*}
%\vec{e}_1\cdot\vec{e}_2
%&=
%\frac{1-2+1}{\sqrt{18}}\\
%&=0\\
%\vec{e}_1\cdot\vec{e}_3
%&=
%\frac{-1+1}{\sqrt{6}}\\
%&=0\\
%\vec{e}_2\cdot\vec{e}_3
%&=
%\frac{-1+1}{\sqrt{12}}\\
%&=0
%\end{align*}
%
%\end{resolution}
%
%%------------------------------------------------------------------------------%
%\begin{resolution}{Verificação}
%\begin{align*}
%\norm{\vec{e}_1}^{2}
%&=
%\frac{3}{3}\\
%&=1\\
%\norm{\vec{e}_2}^{2}
%&=
%\frac{6}{6}\\
%&=1\\
%\norm{\vec{e}_3}^{2}
%&=
%\frac{2}{2}\\
%&=1
%\end{align*}
%\end{resolution}

%------------------------------------------------------------------------------%
\endinput
